% Einheit 1 von 3 – Pyramide (bank/pyramide-kegel-kugel/e1.jsonl, zusammenbau v0.1)
%% TODO Zweigzeile Teil 1: was hier gelernt wird (Satz in Schülersprache aus dem Einheitstitel) – Einheit 1
\einheitenkopf[e1]{Einheit 1 von 3 · Pyramide}
\zweigzeile{ab Kl. 8, je nach Buch bis Kl. 9 · keine P10-Aufgabe}
\setcounter{aufgabe}{11}

%% TODO Titel: Ich-kann-Satz fehlt in der Bank (Platzhalter: Kettenname) – Nr. 12
%% TODO Anweisung: ein Satz über der ersten Teilaufgabe fehlt in der Bank – Nr. 12
\begin{aufgabe}{Säule, Spitze oder Kugel?}
\begin{teile}
\teil Zeltdach – Säule, Spitze oder Kugel? Kreuze an und färbe die Grundfläche.\\ \kreuz{Säule}\\ \kreuz{Spitze}\\ \kreuz{Kugel}
\end{teile}
\end{aufgabe}

%% TODO Titel: Ich-kann-Satz fehlt in der Bank (Platzhalter: Kettenname) – Nr. 13
%% TODO Anweisung: ein Satz über der ersten Teilaufgabe fehlt in der Bank – Nr. 13
\begin{aufgabe}{Stützdreieck einzeichnen}
\swa{Turmdach: Zeichne das rechtwinklige Dreieck aus Höhe, halber Grundkante und Seitenhöhe ein und markiere die Hypotenuse. Nichts rechnen.}{\pyramide{4}{4}{3}{$6$ m}{}{$4$ m}}
\end{aufgabe}

%% TODO Titel: Ich-kann-Satz fehlt in der Bank (Platzhalter: Kettenname) – Nr. 14
%% TODO Anweisung: ein Satz über der ersten Teilaufgabe fehlt in der Bank – Nr. 14
\begin{aufgabe}{Volumen oder Oberfläche?}
\begin{teile}
\teil Stoff für ein pyramidenförmiges Zelt – Volumen oder Oberfläche? Kreuze an.\\ \kreuz{Volumen}\\ \kreuz{Oberfläche}
\end{teile}
\end{aufgabe}

%% TODO Titel: Ich-kann-Satz fehlt in der Bank (Platzhalter: Kettenname) – Nr. 15
%% TODO Anweisung: ein Satz über der ersten Teilaufgabe fehlt in der Bank – Nr. 15
\begin{aufgabe}{Welche Formel?}
\begin{teile}
\teil Pyramide mit der Grundfläche $30$ cm² und der Höhe $8$ cm, gesucht das Volumen – welche Formel? Schreibe sie aus der Formelsammlung ab und die Zahl zu jedem Buchstaben. Nichts rechnen.
\end{teile}
\end{aufgabe}

%% TODO Titel: Ich-kann-Satz fehlt in der Bank (Platzhalter: Kettenname) – Nr. 16
%% TODO Anweisung: ein Satz über der ersten Teilaufgabe fehlt in der Bank – Nr. 16
\begin{aufgabe}{Pyramide}
\begin{teile}
\teil Strecke von der Spitze senkrecht zur Mitte der Grundfläche – Höhe, Seitenhöhe oder Seitenkante? Kreuze an.\\ \kreuz{Höhe}\\ \kreuz{Seitenhöhe}\\ \kreuz{Seitenkante}
\end{teile}
%% TODO Darstellung für schwach fehlt in der Bank (pyramide-kegel-kugel-e1-k5-s1-v1; Abschnitt „Für schwache Schüler“)
\swz{Pyramide: Grundfläche $27$ cm², Höhe $8$ cm – Volumen?}{}
%% TODO Darstellung für schwach fehlt in der Bank (pyramide-kegel-kugel-e1-k5-s1-v2; Abschnitt „Für schwache Schüler“)
\swz{Pyramide: Grundfläche $45$ m², Höhe $6$ m – Volumen?}{}
%% TODO Darstellung für schwach fehlt in der Bank (pyramide-kegel-kugel-e1-k5-s1-v3; Abschnitt „Für schwache Schüler“)
\swz{Pyramide: Grundfläche $21$ dm², Höhe $9$ dm – Volumen?}{}
%% TODO Darstellung für schwach fehlt in der Bank (pyramide-kegel-kugel-e1-k5-s1-v4; Abschnitt „Für schwache Schüler“)
\swz{Pyramide: Grundfläche $64$ cm², Höhe $15$ cm – Volumen?}{}
%% TODO Darstellung für schwach fehlt in der Bank (pyramide-kegel-kugel-e1-k5-s1-v5; Abschnitt „Für schwache Schüler“)
\swz{Pyramide: Grundfläche $36$ m², Höhe $11$ m – Volumen?}{}
\swfrage{Was bleibt gleich, was ändert sich?}
%% TODO Darstellung für schwach fehlt in der Bank (pyramide-kegel-kugel-e1-k5-s2-v1; Abschnitt „Für schwache Schüler“)
\swz{Quadratische Pyramide: Grundkante $6$ cm, Höhe $7$ cm – Volumen?}{}
%% TODO Darstellung für schwach fehlt in der Bank (pyramide-kegel-kugel-e1-k5-s3-v1; Abschnitt „Für schwache Schüler“)
\swz{Pyramide mit rechteckiger Grundfläche $8$ cm mal $5$ cm, Höhe $9$ cm – Volumen?}{}
%% TODO Darstellung für schwach fehlt in der Bank (pyramide-kegel-kugel-e1-k5-s4-v1; Abschnitt „Für schwache Schüler“)
\swz{Quadratische Pyramide: Grundkante $3{,}8$ cm, Höhe $5{,}5$ cm – Volumen? Runde auf eine Stelle.}{}
%% TODO Darstellung für schwach fehlt in der Bank (pyramide-kegel-kugel-e1-k5-s5-v1; Abschnitt „Für schwache Schüler“)
\swz{Quadratische Pyramide: Grundkante $16$ cm, Höhe $15$ cm – Seitenhöhe?}{}
%% TODO Darstellung für schwach fehlt in der Bank (pyramide-kegel-kugel-e1-k5-s6-v1; Abschnitt „Für schwache Schüler“)
\swz{Quadratische Pyramide: Grundkante $8$ cm, Seitenhöhe $11$ cm – Flächeninhalt einer Seitenfläche?}{}
%% TODO Darstellung für schwach fehlt in der Bank (pyramide-kegel-kugel-e1-k5-s7-v1; Abschnitt „Für schwache Schüler“)
\swz{Quadratische Pyramide: Grundkante $6$ cm, Seitenhöhe $9$ cm – Mantel?}{}
%% TODO Darstellung für schwach fehlt in der Bank (pyramide-kegel-kugel-e1-k5-s8-v1; Abschnitt „Für schwache Schüler“)
\swz{Quadratische Pyramide: Grundkante $8$ cm, Seitenhöhe $7$ cm – Oberfläche?}{}
%% TODO Darstellung für schwach fehlt in der Bank (pyramide-kegel-kugel-e1-k5-s9-v1; Abschnitt „Für schwache Schüler“)
\swz{Quadratische Pyramide: Grundkante $6$ cm, Höhe $7$ cm – Seitenkante? Runde auf zwei Stellen.}{}
\swa{Netz einer quadratischen Pyramide (verkleinert): Grundkante $4$ cm, Seitenhöhe $5$ cm. Zeichne in jede Seitenfläche die Seitenhöhe ein und schreibe die Maße an die Grundkanten und Seitenhöhen.}{\netzpyramide{2}{2.5}{}{}}
\swa{Zeichne das Schrägbild einer quadratischen Pyramide mit Grundkante $4$ cm und Höhe $5$ cm (Tiefe halb so lang, schräg unter 45°). Zeichne die Höhe gestrichelt ein und beschrifte sie.}{\rechenplatz{6}}
%% TODO Darstellung für schwach fehlt in der Bank (pyramide-kegel-kugel-e1-k5-s12-v1; Abschnitt „Für schwache Schüler“)
\swz{Pyramide: Volumen $96$ cm³, Grundfläche $36$ cm² – Höhe?}{}
%% TODO Darstellung für schwach fehlt in der Bank (pyramide-kegel-kugel-e1-k5-s13-v1; Abschnitt „Für schwache Schüler“)
\swz{Ein Zeltdach hat die Form einer quadratischen Pyramide ohne Boden: Grundkante $5$ m, Seitenhöhe $3{,}2$ m. Wie viel m² Stoff braucht man?}{}
\swa{Das Dach eines Gartenhauses ist eine quadratische Pyramide mit der Grundkante $4{,}5$ m und der Höhe $1{,}8$ m. Zeichne in das Schrägbild die Höhe des Dachs ein. Beschrifte die Höhe und eine Grundkante mit den Maßen \hfill (P10 2015 OS).}{\pyramide{4}{4}{2}{}{}{}}
\end{aufgabe}

%% TODO Titel: Ich-kann-Satz fehlt in der Bank (Platzhalter: Kettenname) – Nr. 17
%% TODO Anweisung: ein Satz über der ersten Teilaufgabe fehlt in der Bank – Nr. 17
\begin{aufgabe}{Pyramide – Fehler finden, Begründen, Darstellungswechsel, Anwendung}
\swz[3]{Pyramide: Grundfläche $48$ cm², Höhe $9$ cm. Jonas rechnet das Volumen: \rechnung{V &= 48 \cdot 9 \\ V &= 432 \text{ cm}^3} Wo ist der Fehler?}{}
\swfrage{Ein Quader und eine Pyramide haben dieselbe Grundfläche und dieselbe Höhe. Begründe, warum die Pyramide nur ein Drittel so viel Sand fasst.}
\swa{Aus diesem Netz wird eine Pyramide gefaltet. Skizziere ihr Schrägbild und beschrifte Grundkante und Seitenhöhe.}{\netzpyramide{2}{3}{$4$ cm}{$6$ cm}}
\swz[3]{Die Spitze eines Kirchturms ist eine quadratische Pyramide: Grundkante $6$ m, Seitenhöhe $11$ m. Die vier Dachflächen bekommen neues Kupferblech für $85$ € je m². Reichen $11\,000$ €?}{}
\end{aufgabe}

\uebersichtskasten{Volumen: ein Drittel von Grundfläche mal Höhe – drei gleiche Pyramiden füllen das Prisma mit derselben Grundfläche und Höhe. \\ a = 10 cm, h = 12 cm: G = 10 · 10 = 100 cm² \quad V = 100 · 12 : 3 = 400 cm³ \\ Seitenhöhe: rechtwinkliges Dreieck aus der Höhe und der halben Grundkante, die Seitenhöhe ist die Hypotenuse. \quad $h_s$ = √(12² + 5²) = √169 = 13 cm \\ Mantel: vier gleiche Dreiecke, jedes Grundkante mal Seitenhöhe durch zwei. \quad M = 4 · 10 · 13 : 2 = 260 cm² \quad Oberfläche: O = 100 + 260 = 360 cm² \\ Rückwärts: h = 3 · V : G.}
